\documentclass[a4paper,11pt]{article}

% --- Packages ---
\usepackage[utf8]{inputenc}
\usepackage[T1]{fontenc}
\usepackage[french]{babel}
\usepackage{geometry}
\geometry{margin=2.5cm}
\usepackage{graphicx}
\usepackage{xcolor}
\usepackage{amsmath}
\usepackage{amssymb}
\usepackage{booktabs}
\usepackage{array}
\usepackage{longtable}
\usepackage{enumitem}
\usepackage{hyperref}
\usepackage{fancyhdr}
\usepackage{titlesec}
\usepackage{tcolorbox}
\tcbuselibrary{skins, breakable}
\usepackage{colortbl}
\usepackage{multirow}

% --- Colors ---
\definecolor{bmsprimary}{HTML}{1B4F72}
\definecolor{bmssecondary}{HTML}{2E86C1}
\definecolor{bmsgreen}{HTML}{1E8449}
\definecolor{bmsorange}{HTML}{E67E22}
\definecolor{bmsgray}{HTML}{7F8C8D}
\definecolor{bmslight}{HTML}{EBF5FB}
\definecolor{statusdone}{HTML}{1E8449}
\definecolor{statustodo}{HTML}{95A5A6}
\definecolor{bmsred}{HTML}{C0392B}

% --- Header/Footer ---
\pagestyle{fancy}
\fancyhf{}
\fancyhead[L]{\small\textcolor{bmsprimary}{\textbf{Rapport d'Avancement -- BMS Intelligent}}}
\fancyhead[R]{\small\textcolor{bmsgray}{Samsung SDI 60Ah, 1\textsuperscript{er} avril 2026}}
\fancyfoot[C]{\thepage}
\renewcommand{\headrulewidth}{0.4pt}

% --- Title formatting ---
\titleformat{\section}{\Large\bfseries\color{bmsprimary}}{}{0em}{\thesection\quad}
\titleformat{\subsection}{\large\bfseries\color{bmssecondary}}{}{0em}{\thesubsection\quad}
\titleformat{\subsubsection}{\normalsize\bfseries\color{bmsorange}}{}{0em}{\thesubsubsection\quad}

% --- Boxes ---
\newtcolorbox{resultbox}[1][]{
  colback=bmsgreen!5, colframe=bmsgreen,
  fonttitle=\bfseries, title=#1, breakable, left=5pt, right=5pt
}
\newtcolorbox{infobox}[1][]{
  colback=bmslight, colframe=bmsprimary,
  fonttitle=\bfseries, title=#1, breakable
}
\newtcolorbox{warnbox}[1][]{
  colback=bmsorange!8, colframe=bmsorange,
  fonttitle=\bfseries, title=#1, breakable
}

\newcommand{\done}{\textcolor{statusdone}{\textbf{\checkmark}}}
\newcommand{\pending}{\textcolor{statustodo}{\textbf{--}}}

% ===========================================================================
\begin{document}

% --- Title ---
\begin{titlepage}
\centering
\vspace*{1.5cm}
{\LARGE\bfseries\textcolor{bmsprimary}{Rapport d'Avancement}}\\[0.4cm]
{\Large\textcolor{bmssecondary}{Modélisation et Validation d'un BMS Intelligent}}\\[0.3cm]
{\large\textcolor{bmssecondary}{pour Véhicules Électriques -- Samsung SDI 60Ah Prismatic}}\\[1.5cm]

\begin{tcolorbox}[colback=bmslight, colframe=bmsprimary, width=13cm]
\centering
\textbf{Stagiaire}~: Dhiaheddine Kaaniche\\[4pt]
\textbf{Date}~: 1\textsuperscript{er} avril 2026\\[4pt]
\textbf{Avancement}~: Phases 1 et 2 terminées sur 7
\end{tcolorbox}

\vfill
\begin{tcolorbox}[colback=bmsgreen!8, colframe=bmsgreen, width=13cm]
\centering
\textbf{Résumé exécutif}\\[4pt]
Les deux premières phases du projet sont entièrement réalisées~:
(1)~analyse et nettoyage du jeu de données de 70 trajets réels,
(2)~modélisation et calibration du modèle de plant batterie~96S dans
Simulink incluant le modèle ECM~2RC, les tables de paramètres
$R_0/R_1/C_1/R_2/C_2 = f(SoC, T)$, le modèle thermique et la
validation sur l'ensemble des trajets mesurés.
\end{tcolorbox}
\end{titlepage}

\tableofcontents
\newpage

% ===========================================================================
\section{Contexte et Objectifs}
% ===========================================================================

Ce projet de stage vise à concevoir un \textbf{BMS (Battery Management System)
intelligent} pour véhicules électriques, basé sur le jeu de données public
\textit{Battery and Heating Data in Real Driving Cycles} (Samsung SDI 60Ah,
prismatique). L'objectif final est de livrer un système validé selon le
workflow \textbf{MIL/SIL/PIL} conforme à la norme ISO~26262.

Le projet comporte sept phases~: Analyse des données \textrightarrow\
Modélisation Plant \textrightarrow\ Contrôleur BMS \textrightarrow\
Module IA \textrightarrow\ Validation MIL/SIL/PIL \textrightarrow\
Outils IA avancés \textrightarrow\ Intégration finale.

% ===========================================================================
\section{Phase 1 -- Analyse Préliminaire et Exploration des Données}
\textcolor{statusdone}{\textbf{[Terminée]}}
% ===========================================================================

\subsection{Travaux réalisés}

\subsubsection{Contrôle qualité des données brutes}

Les 70 fichiers CSV de mesure (32~TripA~: été, 38~TripB~: hiver) et les
24~fichiers de simulation ont été analysés de manière exhaustive.
Les scripts \texttt{meas\_quality\_check.py} et \texttt{sims\_quality\_check.py}
ont automatisé la détection d'anomalies.

\begin{center}
\small\renewcommand{\arraystretch}{1.3}
\begin{tabular}{lrrr}
\toprule
\textbf{Dataset} & \textbf{Fichiers} & \textbf{Lignes totales} & \textbf{Colonnes} \\
\midrule
Mesure -- TripA (été) & 32 & 467\,701 & 22--28 \\
Mesure -- TripB (hiver) & 38 & 627\,092 & 48 \\
\textbf{Mesure total} & \textbf{70} & \textbf{1\,094\,793} & \textbf{24} \\
\midrule
Simulation (CAN + LIN) & 24 & 917\,281 & 43 \\
\bottomrule
\end{tabular}
\end{center}

Principales anomalies détectées et documentées~:
\begin{itemize}[nosep]
  \item Incohérences de colonnes entre TripA (22--28~colonnes selon le fichier)
    et TripB (48~colonnes avec capteurs thermiques supplémentaires).
  \item Valeurs sentinelles \texttt{65535} dans le signal consommation des fichiers simulation.
  \item En-tête corrompu dans \texttt{Neg10\_Heating\_CAN.csv} (3~lignes d'en-tête au lieu de 2).
  \item Doublons de noms de colonnes (\texttt{Power [kW]} en double) dans les fichiers simulation.
\end{itemize}

\subsubsection{Nettoyage des données}

Les scripts \texttt{meas\_data\_cleaning.py} et \texttt{sims\_data\_cleaning.py}
ont produit deux fichiers consolidés~:

\begin{itemize}[nosep]
  \item \texttt{data/cleaned/measurement/measurement\_clean.csv}~:
    24~colonnes, sans NaN, renommées en \texttt{snake\_case}, identifiant de trip
    (\texttt{trip\_id}) ajouté.
  \item \texttt{data/cleaned/simulation/simulation\_clean.csv}~:
    25~colonnes, sans NaN, identifiant de simulation (\texttt{sim\_id}) ajouté.
\end{itemize}

Traitements appliqués~: suppression des pics capteur par seuillage adaptatif,
interpolation linéaire des lacunes, correction des encodages (latin-1),
sélection des 23~signaux BMS-pertinents.

\subsubsection{Analyse Exploratoire des Données (EDA)}

\begin{center}
\small\renewcommand{\arraystretch}{1.3}
\begin{tabular}{lll}
\toprule
\textbf{Signal} & \textbf{TripA (été)} & \textbf{TripB (hiver)} \\
\midrule
Tension pack [V] & 349 -- 395 & 302 -- 394 \\
Courant batterie [A] & $-$395 à +144 & $-$404 à +145 \\
Température batterie [$^\circ$C] & 16 -- 32 & $-$1 -- 22 \\
Température ambiante [$^\circ$C] & 14 -- 33.5 & $-$3.5 -- 14 \\
SoC [\%] & 34 -- 89 & 0 -- 86 \\
\bottomrule
\end{tabular}
\end{center}

\textbf{Constat important~:} une dérive de +25 à +80~mV par cellule a été
détectée entre TripA et TripB à SoC équivalent, attribuée au vieillissement
ou à la dérive de l'estimateur SoC sur les 5~mois séparant les deux campagnes.
Cette observation impose de traiter les deux campagnes séparément pour
l'extraction des paramètres.

\begin{resultbox}[Livrables Phase 1 produits]
\done~\texttt{ai/reports/data\_quality\_report.md}~-- rapport qualité données brutes\\
\done~\texttt{ai/reports/data\_cleaning\_report.md}~-- rapport nettoyage\\
\done~\texttt{ai/reports/eda\_report.md}~-- analyse exploratoire complète\\
\done~\texttt{data/cleaned/} -- datasets nettoyés prêts à l'emploi\\
\done~Scripts Python~: \texttt{meas\_quality\_check.py}, \texttt{meas\_data\_cleaning.py},
      \texttt{sims\_quality\_check.py}, \texttt{sims\_data\_cleaning.py}
\end{resultbox}

\newpage
% ===========================================================================
\section{Phase 2 -- Modélisation du Plant Batterie}
\textcolor{statusdone}{\textbf{[Terminée]}}
% ===========================================================================

\subsection{Architecture du modèle}

Le modèle de plant est découpé en trois niveaux hiérarchiques dans Simulink~:

\begin{center}
\small\renewcommand{\arraystretch}{1.3}
\begin{tabular}{lll}
\toprule
\textbf{Niveau} & \textbf{Fichier Simulink} & \textbf{Description} \\
\midrule
Cellule & \texttt{cell\_ecm\_2rc.slx} & ECM 2RC avec LUTs $(SoC, T)$ \\
Module & \texttt{module\_12s.slx} & 12 cellules en série \\
Pack & \texttt{pack\_96s.slx} & 8 modules (96S totaux) \\
Refroidissement & \texttt{Coolant\_loop} & Circuit de refroidissement couplé \\
\bottomrule
\end{tabular}
\end{center}

\subsection{Extraction des paramètres ECM 2RC}

Le script Python \texttt{extract\_ecm\_params.py} a extrait les cinq paramètres
du circuit équivalent par optimisation numérique (\texttt{scipy.optimize.least\_squares})
sur les données mesurées~:

\begin{itemize}[nosep]
  \item \textbf{$R_0$}~: identifié par détection des sauts de courant
    ($|\Delta I| > 8$~A/pas).
  \item \textbf{$R_1, C_1, R_2, C_2$}~: ajustement par fenêtres glissantes de
    60~s (stride~30~s, RMSE~$<15$~mV).
  \item \textbf{Grille de LUT}~: $21 \times 7$ points $(SoC~: 0$--$100\%$/5, $T~: 0$--$30\,^\circ$C/5) -- \textbf{données de mesure uniquement}.
  \item Lissage gaussien ($\sigma = [0.8, 0.5]$) pour interpolation stable.
\end{itemize}

\begin{center}
\small\renewcommand{\arraystretch}{1.3}
\begin{tabular}{lcc}
\toprule
\textbf{Paramètre} & \textbf{Plage typique (cellule)} & \textbf{Fichier} \\
\midrule
$R_0$ & 0.4 -- 1.5~m$\Omega$ & \texttt{r0\_lut.csv} \\
$R_1$ & 0.1 -- 1.8~m$\Omega$ & \texttt{r1\_lut.csv} \\
$C_1$ & 3\,000 -- 15\,000~F & \texttt{c1\_lut.csv} \\
$R_2$ & 0.3 -- 5~m$\Omega$ & \texttt{r2\_lut.csv} \\
$C_2$ & 30\,000 -- 150\,000~F & \texttt{c2\_lut.csv} \\
\bottomrule
\end{tabular}
\end{center}

La courbe \textbf{OCV$(SoC)$} a été extraite par le script \texttt{extract\_ocv.py}
et stockée dans \texttt{ocv\_lut.csv}.

\subsection{Modèle du Circuit de Refroidissement (\texttt{Coolant\_loop})}

Un sous-système Simulink \texttt{Coolant\_loop} modélise la dynamique thermique
du liquide de refroidissement~:
\[
  C_{cool}\,\frac{dT_{cool}}{dt}
  = \dot{m}_{cool}\,c_p\,(T_{cool,in} - T_{cool}) + Q_{cell \to cool}
\]
Entrée~: $T_{cool,in}(t)$ issue du signal \texttt{coolant\_temp\_inlet\_c}
(disponible depuis TripA23 et tous les TripB, plage $-1$--$70\,^\circ$C).
Sortie~: $T_{cool}(t)$ injectée dans le modèle thermique de chaque cellule.
La validation a été effectuée sur les trips avec données de refroidissement réelles.

\subsection{Calibration du modèle thermique}

Le modèle thermique par cellule suit l'équation~:
\[
  C_{th}\,\frac{dT_{cell}}{dt}
  = P_{Joule}
  - h_{cool}\,(T_{cell} - T_{cool})
  - h_{amb}\,(T_{cell} - T_{amb})
\]

La calibration a été effectuée par évaluation exhaustive de candidats (script
\texttt{calibrate\_thermal.m}) sur 9~trips représentatifs (TripA23+/TripB,
température refroidissant disponible).

\begin{center}
\small\renewcommand{\arraystretch}{1.3}
\begin{tabular}{lrl}
\toprule
\textbf{Paramètre} & \textbf{Valeur calibrée} & \textbf{Signification} \\
\midrule
$C_{th}$ & 1300~J/K & Capacité thermique cellule \\
$h_{cool}$ & 0.02~W/K & Couplage cellule--liquide de refroidissement \\
$h_{amb}$ & 0.15~W/K & Échange cellule--ambiant \\
\midrule
RMSE thermique & 0.75\,$^\circ$C & Sur 9 trips de calibration \\
\bottomrule
\end{tabular}
\end{center}

\begin{warnbox}[Observations importantes]
\begin{itemize}[nosep]
  \item TripA01--A22~: température du liquide de refroidissement indisponible
    (tous zéros) -- ces trips n'ont pas été utilisés pour la calibration thermique.
  \item Le liquide de refroidissement affiche 30--60\,$^\circ$C de plus que les cellules
    en conditions hivernales (chauffage actif du TMS).
  \item Un bug d'initialisation a été corrigé~: \texttt{T\_cells\_init} utilisait
    par erreur $T_{cool}(t=0)$ au lieu de $T_{cell}(t=0)$, provoquant une erreur
    thermique initiale de 50\,$^\circ$C.
\end{itemize}
\end{warnbox}

\subsection{Résultats de validation du pack 96S}

Le script \texttt{validate\_pack\_model.m} a simulé l'ensemble des 70~trips et
comparé la tension pack et la température moyennes modélisées aux mesures.

\begin{center}
\small\renewcommand{\arraystretch}{1.4}
\begin{tabular}{lcccc}
\toprule
\textbf{Indicateur} & \textbf{Médiane} & \textbf{Moyenne} & \textbf{Max} & \textbf{Unité} \\
\midrule
RMSE tension pack & 1\,102 & 1\,563 & 6\,201 & mV \\
Biais tension & -- & $-591$ & -- & mV \\
RMSE température & 1.18 & 1.46 & 4.59 & $^\circ$C \\
RMSE SoC & 0.18 & 0.25 & 0.65 & \% \\
\bottomrule
\end{tabular}
\end{center}

\begin{infobox}[Domaine de validité du modèle]
Le modèle ECM~2RC du pack~96S est \textbf{validé dans l'enveloppe opérationnelle
suivante}, directement issue de la couverture du dataset~:
\begin{itemize}[nosep]
  \item $SoC \in [40\%, 90\%]$~: plage effectivement couverte par les données de mesure.
  \item $T_{cell} \in [10\,^\circ C, 30\,^\circ C]$~: plage à forte densité de points.
\end{itemize}
Dans ce domaine, la RMSE~V est de l'ordre de 200--700~mV (pack), soit
$\approx 2$--$7$~mV par cellule. Les conditions hors-enveloppe ($SoC < 40\%$,
$T < 10\,^\circ$C) présentent des erreurs plus élevées (2\,500--6\,000~mV),
cohérent avec l'absence de données de caractérisation à ces niveaux.
Les données de simulation (3~températures~: $-10$, $0$, $+10\,^\circ$C)
ne permettent pas de combler les bins manquants.
\end{infobox}

\begin{resultbox}[Livrables Phase 2 produits]
\done~\texttt{cell\_ecm\_2rc.slx}, \texttt{module\_12s.slx}, \texttt{pack\_96s.slx} -- modèles Simulink\\
\done~Sous-système \texttt{Coolant\_loop} (circuit de refroidissement) intégré\\
\done~\texttt{plant\_model/data/\{r0,r1,c1,r2,c2\}\_lut.csv} -- tables LUT $21\times7$ (mesure)\\
\done~\texttt{plant\_model/data/ocv\_lut.csv} -- courbe OCV$(SoC)$\\
\done~\texttt{plant\_model/data/ecm\_params.mat} -- paramètres MATLAB\\
\done~\texttt{plant\_model/data/validation\_results.csv} -- métriques des 70~trips\\
\done~\texttt{plant\_model/data/drivecycles\_meas.mat} -- profils de conduite (mesure)\\
\done~Scripts~: \texttt{extract\_ecm\_params.py}, \texttt{extract\_ocv.py},
      \texttt{calibrate\_thermal.m}, \texttt{validate\_pack\_model.m},
      \texttt{init\_plant\_model.m}, \texttt{export\_luts.m},
      \texttt{export\_drivecycles.m}
\end{resultbox}

\newpage
% ===========================================================================
\section{Tableau de Bord Global}
% ===========================================================================

\begin{center}
\renewcommand{\arraystretch}{1.5}
\begin{tabular}{clp{5.5cm}c}
\toprule
\textbf{Phase} & \textbf{Intitulé} & \textbf{Avancement} & \textbf{Statut} \\
\midrule
1 & Analyse Préliminaire
  & EDA, qualité, nettoyage terminés
  & \textcolor{statusdone}{\textbf{\checkmark~Terminé}} \\
2 & Modélisation Plant
  & ECM 2RC + thermique + validation
  & \textcolor{statusdone}{\textbf{\checkmark~Terminé}} \\
3 & Contrôleur BMS
  & EKF SoC, SoH, Balancing, Fault Mgr
  & \textcolor{statustodo}{\textbf{À faire}} \\
4 & Module IA
  & Fault injection, LSTM, Simulink
  & \textcolor{statustodo}{\textbf{À faire}} \\
5 & Validation MIL/SIL/PIL
  & Génération code C, PIL MCU
  & \textcolor{statustodo}{\textbf{À faire}} \\
6 & Outils IA Avancés
  & Dead code, complexité, doc auto
  & \textcolor{statustodo}{\textbf{À faire}} \\
7 & Intégration Finale
  & Assemblage, ISO 26262, rapport
  & \textcolor{statustodo}{\textbf{À faire}} \\
\bottomrule
\end{tabular}
\end{center}

% ===========================================================================
\section{Prochaines Étapes (Phase 3)}
% ===========================================================================

\begin{enumerate}[nosep]
  \item \textbf{Implémenter l'estimateur SoC EKF}~:
    état $[SoC, V_{RC1}, V_{RC2}]^T$, Jacobiens, réglage $\mathbf{Q}/R$,
    objectif RMSE~$<2\%$.
  \item \textbf{Implémenter l'estimateur SoH}~:
    tracking de $R_0$ par RLS, estimation de capacité résiduelle.
  \item \textbf{Développer le Stateflow \texttt{Fault\_Manager}}~:
    surtension, sous-tension, surcourant, surchauffe, défaut capteur~;
    logique de debounce et DTC ISO~26262.
  \item \textbf{Développer le Stateflow \texttt{Cell\_Balancing}}~:
    équilibrage passif sur les 96 cellules du pack.
\end{enumerate}

% ===========================================================================
\section{Points d'Attention}
% ===========================================================================

\begin{warnbox}[Points à prendre en compte pour la Phase 3]
\begin{enumerate}[nosep]
  \item \textbf{Domaine de validité du plant~:} le modèle est validé pour
    $SoC \in [40\%, 90\%]$ et $T \in [10, 30]\,^\circ$C. L'EKF devra être
    testé et réglé dans cette enveloppe.
  \item \textbf{Dérive TripA/TripB~:} prendre en compte la dérive de
    vieillissement (+25--80~mV/cellule) lors de la validation de l'EKF --
    initialiser les LUTs sur TripB séparément.
  \item \textbf{Couverture thermique réduite~:} les TripA01--A22 n'ont pas de
    données de refroidissement~; le \texttt{Coolant\_loop} fonctionnera en
    mode refroidissement naturel ($h_{cool} = 0$) pour ces trips.
\end{enumerate}
\end{warnbox}

\end{document}
